\section{A stronger initialization consequence: the mean from every safe start}
\label{sec:anystart}
The mean-only decision does not require acquiring or preparing an entire
equilibrium initial law in the certified boundary regime. The same finite
certificate controls every safe starting configuration. This changes the
initialization operation while keeping the original mean objective as its
benchmark. It does not supply the earlier marked-law or CDF initial law.

Use the true-generator kernel $P^\circ$, and the $k,L,z,a,g_0,\bar p$
of Section~\ref{sec:boundary}. Define
\begin{align}
 s_*&=(1-\bar p)^2,\qquad
 E_* = La^k+a^{k/2}e^{-g_0L}/g_0,\\
 A_*&=(1-z)(1-\bar p)(1-E_*/s_*),\\
 B_*&=1+(L\bar p+e^{-g_0L}/g_0)/s_*.
 \label{eq:safestartconstants}
\end{align}
Require $E_*<s_*$. All these are certificate arithmetic; $L$ is not an
operating waiting time.

\begin{theorem}[Uniform safe-start mean comparison]
For every menu item and every safe state $x\ne b$,
\begin{equation}
 A_* M_u\le\E_{x,Q^u}\tau_b\le B_* M_u.
 \label{eq:allsafemean}
\end{equation}
The upper bound also holds at $b$, taking $\tau_b=0$. Consequently any
actual starting law $\nu_u$ with $\nu_u(b)\le f$ satisfies
\begin{equation}
 (1-f)A_* M_u\le\E_{\nu_u,Q^u}\tau_b\le B_*M_u.
 \label{eq:imperfectsafemean}
\end{equation}
These are bounds for an unbounded mean, not deductions from a TV estimate.
\end{theorem}
\begin{proof}
For $Z_u(x)=\sum_{j\ge0}[(P^\circ)^j(x,b)-p]$ and $S_u=Z_u(b)$,
\[
 \frac{\E_{x,Q^u}\tau_b}{M_u}
 =(1-p)\left(1-\frac{Z_u(x)}{S_u}\right),\qquad S_u\ge s_*.
\]
At any $x\in C_k$, the finite sum and spectral bounds give
$Z_u(x)\le E_*$. Thus the ratio there is at least
$(1-\bar p)(1-E_*/s_*)$.
For a safe $x$ with $w(x)<k$, the same decreasing superharmonic
function $H$ gives
$\Pr_x(\tau_b<\tau_{C_k})\le H(w(x))\le H(1)\le z$.
Apply the strong Markov property at the first hit of $\{b\}\cup C_k$;
the time already spent is nonnegative. This proves the lower bound.

For every state, $\pi(b)/\pi(x)\le1$. Each of the first $L$ terms of
$Z_u(x)$ is at least $-p$, and the remaining spectral sum has modulus
at most $e^{-g_0L}/g_0$. Hence
$Z_u(x)\ge-Lp-e^{-g_0L}/g_0$, giving the upper bound.
Integrate over $\nu_u$ and use the zero lifetime assigned to an
initially failed state.
\end{proof}

At fixed local margins, choosing the finite tails sufficiently small
makes $A_*/B_*$ approach one as $N$ grows, with the explicit
$z=O(1/N)$ return correction. Therefore this is not an assertion that
all safe starts are equivalent in arbitrary rare-event systems. It is
a consequence of the uniform recovery drift, singleton target and
mixing certificate in this family.

\subsection{The preferred mean-only preparation: a bounded operating wait}
An additional graphical comparison removes inspection when a small
explicit initial-failure mass is allowed by the mean objective. Put
$a_{\rm w}=r_-(1-\bar q)$ and $c_{\rm w}=r_+\bar q$.
Couple the actual failed-bit process $X$ to independent bits $Y_i$,
initially all failed, with recovery rate $a_{\rm w}$ and failure rate
$c_{\rm w}$. A rate-$c_{\rm w}$ clock sets $Y_i$ to failed and gives
$X_i$ its appropriately thinned failure transition. A rate-$a_{\rm w}$
clock recovers both bits, and extra recovery clocks supply the remaining
actual recovery rate. Every actual initial state obeys $X\le Y$, and
these clocks preserve that order. This comparison does not require the
actual interacting system to be attractive; the clocks are a proof
device, not additional physical control.

Thus, from any initial law, at deterministic operating time $t$,
\begin{equation}
 \Pr(X_t=b)\le q_{\rm dom}(t)^N,\qquad
 q_{\rm dom}(t)=\frac{c_{\rm w}}{a_{\rm w}+c_{\rm w}}
 +\frac{a_{\rm w}}{a_{\rm w}+c_{\rm w}}
                  e^{-(a_{\rm w}+c_{\rm w})t}.
 \label{eq:dominatedwarm}
\end{equation}
Every $(N-1)$-coordinate marginal satisfies the analogous bound.
For an actual terminal time in $[D_-,D_+]$, put
$q=q_{\rm dom}(D_-)$ and $\Delta_D=D_+-D_-$.
A visit to $b$ in that window requires either already being there or
an entry caused by one component's failure. Its expected entry count is
at most $Nc_{\rm w}\Delta_D q^{N-1}$. Therefore
\begin{equation}
 \Pr(\exists t\in[D_-,D_+]:X_t=b)
 \le\chi_D:=\min\{1,q^N+Nc_{\rm w}\Delta_Dq^{N-1}\}.
 \label{eq:finitewarm}
\end{equation}
The inclusion of the actual terminal event in this entire-window event
also permits bounded adaptive terminal timing. The actual operating
dynamics must remain in force during commissioning; no field-biased
clock is substituted. The timer must supply a finite upper and lower
elapsed-time bound, and its actual response/calibration cost is paid.

Install and settle the selected intervention, run this finite warm-up,
then begin the mission unconditionally. The actual initial law has
$\nu_u(b)\le\chi_D$, so~\eqref{eq:imperfectsafemean} applies directly.
There is no inspection, rejection, abort or equilibrium-law claim in
this branch. An initially failed mission has lifetime zero. If
$D_-=\log2/(a_{\rm w}+c_{\rm w})$ and the width $\Delta_D$ is fixed,
$\chi_D$ decreases exponentially in $N$, with the explicit linear
prefactor. The physical commissioning time is bounded independently
of $N$ at fixed local margins, while summed component exposure is
$ND_+$; installation, propagation, finite slew, settling and timer
services remain additional. Repairability during commissioning remains
essential. If the declared target loss cannot accommodate $\chi_D$,
use the next interface or the stronger paid relative preparation.

\subsection{A finite one-bit alternative at the actual operating setting}
This physical realization requires a repairable installation and a
noninvasive meter for one chosen actual component. The earlier full
safe-state inspection access supplies such a meter; local copies alone
would not. The meter must capture that component's binary state at a
declared sample time and report by a finite deadline. A finite capture-time
uncertainty is included in the elapsed-delay allowance below. This is
stated access, not a newly invented sensor. Its response, calibration,
latency, setup and read-error prices remain explicit. The capture time and
its finite jitter are exogenous to the current focal path conditional on
previous history; a path-adaptive capture needs an additional error bound.

Let the system evolve at the selected actual menu from any starting
configuration. For the chosen bit, put
\[
 a_{\rm w}=r_-(1-\bar q),\qquad c_{\rm w}=r_+\bar q.
\]
Its instantaneous failed-to-working rate is at least $a_{\rm w}$ and
its working-to-failed rate is at most $c_{\rm w}$, for every neighborhood.
If $v(t)$ is its working probability, then
\[
 v'(t)\ge a_{\rm w}(1-v(t))-c_{\rm w}v(t).
\]
No independence of the components is used. Thus an exogenous warm-up
of duration at least $t_0$ gives
\begin{equation}
 v(t)\ge v_0:=\frac{a_{\rm w}}{a_{\rm w}+c_{\rm w}}
       \left(1-e^{-(a_{\rm w}+c_{\rm w})t_0}\right).
 \label{eq:witnesswarm}
\end{equation}
The bound is uniform over the entire history at the start of an attempt.
A convenient finite choice is
$t_0=\log2/(a_{\rm w}+c_{\rm w})$, with an outward-rounded timer.

Suppose the captured-bit error is at most $\beta_{\rm read}$ conditional
on its true state and previous history, and the mission begins no later
than $\Delta_{\rm read}$ after the earliest possible capture time. The
operating dynamics continue during this interval. Accept only a reported
working bit. Then
\begin{equation}
 s_{\rm acc}\ge(1-\beta_{\rm read})v_0,\qquad
 f_{\rm start}\le
 \frac{\beta_{\rm read}+c_{\rm w}\Delta_{\rm read}}
      {(1-\beta_{\rm read})v_0}.
 \label{eq:witnesserrors}
\end{equation}
For the second inequality, an accepted but globally failed terminal
state requires either a wrong captured-bit report or a subsequent failure
of that bit. A correct working bit has probability at most
$c_{\rm w}\Delta_{\rm read}$ of such a transition. Dividing by the
acceptance lower bound proves the claim. An erroneous acceptance is
therefore retained as possible initial failure, not relabeled a certainly
safe preparation. Detector common-mode failures must be charged in the
common calibration confidence or the stated conditional bound.

Repeat up to
\[
 J_{\rm prep}=\left\lceil
 \frac{\log\delta_{\rm abort}}{\log(1-s_{\rm acc})}
 \right\rceil
\]
attempts; return an explicit abort if none accepts. Uniform conditional
bounds give abort probability at most $\delta_{\rm abort}$ even for
dependent reused apparatus. The mixture over accepted attempts still
has $\nu_u(b)\le f_{\rm start}$. No equilibrium sampler, rare-state
rejection, complete state readout or full-system pin is required.

Each attempt pays its actual whole-system warm-up and latency time,
one capture/read, controller/timer response, and any reset/setup.
Summed component exposure is at most
$N J_{\rm prep}(t_{\rm warm}^{+}+\Delta_{\rm read})$, with other active
installation intervals added. At fixed local bounds and fixed desired
$f_{\rm start},\delta_{\rm abort}$ this is $O(N)$; it is not constant
total exposure merely because the installation evolves in parallel.
Field installation, propagation, slew, settling and the stated meter's
physical price are additional. If the source lacks the required one-bit
meter, or its latency/read fidelity cannot satisfy~\eqref{eq:witnesserrors},
this realization is unavailable and the separately paid preparation must
be used. The process must be repairable before mission start.

\subsection{The complete deployed decision factor}
Let $F_{\rm dec}$ denote the proved selected-to-best ratio for the
equilibrium-safe benchmark, for example
$(\ell/U)(1-\eta)/(1+\eta)$ in~\eqref{eq:groupedglobal}.
For the just-constructed preparations, the actual means
$\mathcal M_u=\E_{\nu_u,Q^u}\tau_b$ obey
\begin{equation}
 \boxed{\frac{\mathcal M_{\widehat u}}{\max_u\mathcal M_u}
 \ge (1-f_{\rm start})\frac{A_*}{B_*}F_{\rm dec}.}
 \label{eq:deployedsafe}
\end{equation}
For the unconditional warm-up branch take $f_{\rm start}=\chi_D$ and
there is no preparation abort. For the one-bit branch and a mission
definition assigning lifetime zero to a preparation abort,
multiply the lower guarantee by $1-\delta_{\rm abort}$. The confidence
of the acquisition/decision certificate remains a separate probability.
All displayed factors are checked against the requested tolerance before
issuing the decision. Additional deployed-generator error is allowed only
under its separately proved common-law form bounds or by widening the
actual operating certificates; it is not hidden inside detector error.

This completes the mean-only initialization interface. Relative equilibrium
preparation remains a distinct stronger operation for outputs that require
that initial law. The next section retains its conservative bill; the
independent logarithmic-Sobolev construction in the archive improves its
mixing schedule under its additional acquired norm condition.
